\appendix

\renewcommand{\thesubsection}{S\arabic{subsection}}
\renewcommand{\theequation}{S\arabic{equation}}
\renewcommand{\thefigure}{S\arabic{figure}}
\setcounter{figure}{0}
\counterwithin*{equation}{section}

\section*{Supplementary Information}

\textbf{\huge{Tunable Synaptic Working Memory with Volatile Memristive Devices}}\\ \\

\author{Saverio Ricci*$^{1}$, David Kappel*$^{2,3}$, Christian Tetzlaff$^{2,4}$, Daniele Ielmini$^{1}$, and Erika Covi$^{5}$\\
\\
$^{1}$ Dipartimento di Elettronica, Informazione e Bioingegneria (DEIB), Politecnico di Milano, Piazza Leonardo da Vinci 32, 20133 Milano, Italy \\
$^{2}$ Bernstein Center for Computational Neuroscience,  III Physikalisches Institut – Biophysik, Georg-August Universität, Friedrich-Hund-Platz 1, 37077 Göttingen, Germany \\
$^{3}$ Institut für Neuroinformatik, Ruhr-Universität Bochum, Universitätsstr. 150 NB 3/32,
44801 Bochum, Germany \\
$^{4}$ Group of Computational Synaptic Physiology, Department for Neuro- and Sensory Physiology, University Medical Center G\"ottingen, Von-Siebold-Str. 3A, 37075 G\"ottingen, Germany \\
$^{5}$ NaMLab gGmbH, N\"othnitzer Str. 64 a, 01187 Dresden, Germany \\
* These Authors contributed equally to this work.}


%\newpage
\begin{figure}[ht]
    \centering
    \includegraphics[width=0.95\textwidth]{Figures/QuasiStaticCharacterization.png}
    \caption{Quasi-static characterization of the devices. a) First cycle (dashed lines) and subsequent cycles at different compliance currents.  b) Comparison of 5 different devices.} %The devices have the same behavior and repeatable shape.}
    \label{fig:Supfig_01}
\end{figure}

Fig.~\ref{fig:Supfig_01}a reports an analysis of quasi-static measurements conducted on an Ag-based volatile memory. %The I-V sweeps are carried out using Agilent HP4156C semiconductor parameter analyzer, the maximum amplitude is 1.5V and variable compliances are applied. 
The characterization highlights that the forming (dashed line), or soft breakdown of the dielectric layer, has a voltage equal to the set ones and the set voltages are independent on the compliance currents. %The forming initialization, soft breakdown of the dielectric layer, is highlighted with a dashed black line and it is similar to the others I-V curves. The compliance current seems to not affect the I-V response of the device, in good agreement with the volatile behavior of such RRAMs. The device, in fact, cannot know in advance the compliance value and thus this parameter cannot affect the set transition. 

The conductance of the device, measured at 0.1\,V, passes from a \gls{hrs} $\ge$ 10\,T$\Omega$ to a \gls{lrs} $\approx$\,10\,k$\Omega$, with linear ohmic behavior in this last condition. The R$\mathrm{_{ON}}$ of the MOS transistor is in the order of a few hundreds ohm, well below the device resistance.

Fig.~\ref{fig:Supfig_01}b compares the I-V curves of the 5 devices used for the experiments, which show small device-to-device variability.

\newpage


\begin{figure}[ht]
    \centering
    \includegraphics[width=0.95\textwidth]{Figures/MultipleSwithingMeasurement.png}
    \caption{%RRAM responses to a group of pulses with different switching probabilities. 
    Switching of volatile memristive devices when stimulated with different number of pulses and voltage amplitudes. The trains of pulses are applied at 1\,s time interval. The presented images are cropped (as shown on top).}
    \label{fig:Supfig_02}
\end{figure}


To study the switching probability for a given voltage as a function of the number of pulses (P$_{ON}$(N)), trains of N 1\,ms pulses with a period of 11\,ms are applied. Each train is preceded by a 10\,ms reading pulse to verify that the device is initially in its \gls{hrs}. During this characterization, the I$_C$ was 17\,\textmu A, resulting in an average retention time of 28\,ms. Therefore, the trains of pulses are applied at 1\,s time intervals to allow the device to switch off. Fig.~\ref{fig:Supfig_02} compares the impact of the number of pulses (N) per train.

The smaller the voltage amplitude, the smaller the switching probability (P$_{ON}$(V)). For N = 6, for example, the device switches ON after the 4$^{th}$ pulse of the first train, whereas during the 3$^{rd}$ train the device does not switch ON. Increasing the number of pulses in the train to N = 20 increases the probability of the device to be in \gls{lrs} at the end of the train.

\newpage

\begin{figure}[ht]
    \centering
    \includegraphics[width=0.95\textwidth]{Figures/SupplementaryTraining.png}
    \caption{Probability of switching a device ON in a certain time window as a function of the stimulation rate and the switching probability.}% The longer the time window, the higher the probability of the device to be in the \gls{lrs}.}
    \label{fig:Supfig_03bis}
\end{figure}

Fig.~\ref{fig:Supfig_03bis} compares different time windows in which a volatile device is continuously stimulated with a constant stimulation rate for a given switching probability.

Within the same time window, the probability of switching a device ON increases increasing either the switching probability (i.e., the voltage amplitude) or the stimulation rate, which is linked to the number of pulses in the train. As an example, with a time window of 200\,ms, the probability to find a device ON after stimulation ranges from almost 0\% with a stimulation rate of 10\,Hz and switching probability of 5\% to about 80\% with a stimulation rate of 50\,Hz and switching probability of 20\%.

If the stimulation rate and switching probability are fixed, longer time windows increase the probability of the device to be ON at the end of the stimulation sequence. If we consider a 10\,Hz stimulation with 5\% switching probability, the probability to find a device ON after stimulation ranges from almost 0\% with a time window of 100\,ms to 100\% with a time window of 5\,s. %For short time windows, and thus a small number of pulses, the probability of the device to switches to the \gls{lrs} is very low. As an example, in the 100\,ms window only 1 (10\,Hz), 2 (20\,Hz), and 5 (5\,Hz) pulses are applied. By considering longer monitoring windows, the probability increases according to the larger number of pulses the device receives.

\newpage

\begin{figure}[ht]
    \centering
    \includegraphics[width=0.95\textwidth]{Figures/WorkingMemoryTraceExample.png}
    \caption{\gls{wm} store and recall experiment. a) The top panel shows the color input stream. Each color is encoded in a specific pattern (middle panel)%, in that case the green color is the 01101 pattern. The 1 is encoded in a voltage pulse while the 0 is encoded in a zero pulse of the same duration. 
    The voltage amplitudes are chosen to have a switching probability of 5\%. The bottom panel shows the current collected from the devices that flows into the neuron. During the store phase the 3 devices that encode the color green switch on at different times. During the recall phase, the color is recognized if the current from the memristive synapses exceeds a threshold of 43\,\textmu A. In the forgetting phase, the 3 devices switches off at different times, thus indicating that the network forgot the stored element. b) Zoom of a part of the recall phase. Only the stored pattern, in the blue box, provides a current above the threshold, thus triggering the firing of the neuron.} %the three switches of the stimulated devices during the training. A compliance current of 17\,\textmu A is selected, allowing discretized levels. During the recall random patterns are sent to the system. Depending on the number of ON cells stimulated a precise current level is reached and a threshold trigger the system. After the experiment a small reading voltage is applied to monitor the retention capability. The retention current monitor the switching OFF of the devices. b) Insight of a recall part. The 01101 pattern was stored. The right patters, in the blue box, stimulate all the three ON devices and 3 levels of current is flowing (53\,\textmu A).  When only two ON devices are stimulated there are two levels (37\,\textmu A) and when only one is stimulated the current reaches 17\,\textmu A. A current threshold of 43\,\textmu A is used to distinguish when the right patter feds the system from the other patterns.}
    \label{fig:Supfig_03}
\end{figure}

Fig.~\ref{fig:Supfig_03} reports the example of an experimental trace, presenting the store, recall, and forget phases. %by looking at the patterns (encoded with the 10 different colors), the voltages across the RRAMs and the overall system current response.
During the store phase, a single pattern (the green one in the example, encoded as 01101) is applied multiple times and until the three devices associated with 1 are in \gls{lrs}. The voltages applied to the devices and shown in the middle panel of Fig.~\ref{fig:Supfig_03}a are chosen to have a 5\% switching probability (P$_{ON}$). As a consequence, the total current of the system increases (bottom panel). Since the devices have all the same current compliance of 17\,\textmu A, each device that switches on provides a contribution in current of 17\,\textmu A.

During the recall phase, random patterns are fed to the system. The current contribution of each stimulated device in \gls{lrs} is 17\,\textmu A, therefore we expect 51\,\textmu A when 3 ON devices are stimulated (stored pattern) and 34\,\textmu A when 2 ON devices are stimulated (worst case for the non-stored patterns). Therefore, a threshold of 42\,\textmu A is set. When the stored pattern is presented to the network, the current contribution exceeds the threshold, thus triggering the fire of the neuron. The current contribution of the other patterns is lower than the threshold, as better highlighted in Fig.~\ref{fig:Supfig_03}b.

After the recall phase, the system is monitored in absence of stimulation with a reading voltage of -150\,mV applied to all the devices. After about 1\,s, all the devices switch off, thus indicating that the network forgot the stored element. As in the store phase the devices have different switching on times, also the switch off occurs with different times due to the stochastic nature of the \gls{cf}.
%are random and the system responds as a consequence of the number of ON devices are stimulated. When the true pattern arrives (the green one in the example) the system has three levels of curremt, since all the 3 ON devices are stimulated. When only two ON RRAMs are stimulated the current reaches the second level, and a similar reasoning takes place when only 1 ON RRAM is stimulated. Fig.S3b better highlights this concept, comparing the current result for the different patterns. At the end of the experiment, -100 mV bias voltage is used to monitor the system status and the temporal evolution of the architecture.


\newpage

\begin{figure}[ht]
    \centering
    \includegraphics[width=0.8\textwidth]{Figures/LevelCurrentDistributions.png}
    \caption{Currents distribution of the 5 devices used in the store and recall experiment. a) Current read when all devices are in \gls{hrs} and when one, two, or three are in \gls{lrs}. Data grouped by stored pattern. The dashed lines separate number of ON devices. %Threshold current of 42\,\textmu A is a suitable solution for the reservoir system. 
    b) Current distribution of the currents from all the patterns combined. 42\,\textmu A threshold is set as threshold.}
    \label{fig:Supfig_04}
\end{figure}

A threshold current is used to recognize the stored pattern from the other. According to the current levels presented in Fig.~\ref{fig:Supfig_04}, a 42\,\textmu A threshold is suitable to distinguish the stored pattern, since there is a clear non-overlap among the different current distributions.

\newpage

\begin{figure}[ht]
    \centering
    \includegraphics[width=0.95\textwidth]{Figures/LearingComparison.png}
    \caption{Impact of the switching probability on the time needed to store a pattern. The stimulation rate is with 20\,Hz, whereas the switching probability is changed from 1\% to 30\%. The blue numbers inside each plot indicate how many pulses are necessary before a device is set to \gls{lrs}.}
    \label{fig:Supfig_05}
\end{figure}
 
Fig.~\ref{fig:Supfig_05} shows the effect of the switching probability on the store phase. Low switching probabilities require long trains of pulses to set the stimulated devices in \gls{lrs}. As an example, with P$\mathrm{_{ON}}$\,=\,1\%, 10\,s stimulation are not enough to store a pattern and only one devices is in \gls{lrs} by the end of the stimulation. An increase of P$\mathrm{_{ON}}$ shortens the time needed to store a pattern, which can be don to less than 1\,s in the case of P$\mathrm{_{ON}}$\,=\,30\%. It should be noted, however, that a high P$\mathrm{_{ON}}$ might be beneficial in storing a pattern in a short time, but in the recall phase leads to an early forgetting of the pattern. It is therefore important to select the P$\mathrm{_{ON}}$ while considering all the phases of the \gls{wm} system.

\newpage

\begin{figure}[ht]
    \centering
    \includegraphics[width=0.95\textwidth]{Figures/LevelsMismatch.png}
    \caption{Failure analysis during the recall phase of the \gls{wm} store and recall experiment. a) Example of failure defined as the mismatch between the number of actual ON devices and the number of expected ones when a pattern is presented to the network. b-f) Single plots of the data in a) to highlight the behavior of the \gls{wm} system under different P$\mathrm{_{ON}}$ and stimulation rates.}%Errors can be a ON device which switches OFF or an OFF device which switches ON. When an error occurs (ON or OFF cell) the impact of that error is visible in the trace. Despite the error the system has a kind of robustness and zero mismatch is still possible. b) When the reservoir behaves as expected, the mismatch is zero. c) 2 wrong devices switch ON. d) 3 right devices switch OFF. e) 1 cell switches OFF and 2 cells switch ON. f) 1 cell switches ON.}
    \label{fig:Supfig_06}
\end{figure}

The possible failures of the system can be measured by comparing the expected and read currents, which correspond to the expected and actual number of devices in \gls{lrs} contributing to the read current (Fig.~\ref{fig:Supfig_06}a). Four possible situations can arise, summarized in Fig.~\ref{fig:Supfig_06}:
\begin{itemize}
    \item No difference between expected and read current (Fig.~\ref{fig:Supfig_06}b), no failures.
    \item Expected current lower than read current (Fig.~\ref{fig:Supfig_06}c and f). This is usually the case when the switching probability and/or the stimulation rate are too high. One or more devices, previously in \gls{hrs}, are switched to \gls{lrs} during the recall phase.
    \item Expected current greater than read current (Fig.~\ref{fig:Supfig_06}d). In this case case the switching probability and/or the stimulation rate are too low and one or more devices, previously in \gls{lrs}, switch back to \gls{hrs} during the recall phase.
    \item Combination of the two previous cases (Fig.~\ref{fig:Supfig_06}e). This is an unstable situation in which during the recall phase some devices switch from \gls{lrs} to \gls{hrs} and vice versa. In this case, the stochastic nature of the memristive devices is the main reason of the fluctuations observed.
\end{itemize}

\newpage

\begin{figure}[ht]
    \centering
    \includegraphics[width=0.95\textwidth]{Figures/CorrelationPlots.png}
    \caption{Correlation plots of the system response. a) The measured current is shown on the horizontal axis and the expected current is on the vertical axis. The color-bar, as well as the numbers inside each square indicate the probability of each occurrence. b) Correlation plots for different experimental conditions.}
    \label{fig:Supfig_07}
\end{figure}

Fig.~\ref{fig:Supfig_07} shows the correlation plot of different experiments carried out by varying the switching probability and the repetition rate. Lower P$\mathrm{_{ON}}$ improves the classification accuracy of the system. When the stimulation rate is lower than the average retention time of the device, the devices tend to switch from \gls{lrs} to \gls{hrs} during the recall phase, thus leading to higher probability of forgetting the pattern. High stimulation rates instead, cause devices to switch from \gls{hrs} to \gls{lrs}, thus overwriting the stored pattern. A careful choice of average retention times and stimulation rate has therefore to be made to maximize the system accuracy.

\newpage

\begin{figure}[ht]
    \centering
    \includegraphics[width=0.95\textwidth]{Figures/working-memory-network-detail.png}
    \caption{Details to the working memory network model. Connection probabilities (in \%) and PSP amplitude shown between neuronal populations. Input and working memory populations only shown for two memory items to keep the figure uncluttered.}
    \label{fig:Supfig_08}
\end{figure}

\newpage

\begin{figure}[ht]
    \centering
    \includegraphics[width=0.95\textwidth]{Figures/lognormal-fit.png}
    \caption{Lognormal fits (solid lines) to empirical cumulative distribution of retention times for different currents (data points).}
    \label{fig:Supfig_lognormal_fit}
\end{figure}

Fig.~\ref{fig:Supfig_lognormal_fit} shows the lognormal fit to the experimental data of the memristive device. The device is stimulated with semi-triangular pulses as in Fig.~\ref{fig:fig2} and the current compliance was varied by changing the voltage at the gate of the transistor.